\section{Problems we hit, and claims we withdrew}\label{sec:problems}

We list these because a reader who builds on this work will otherwise meet them again.

\begin{enumerate}[leftmargin=1.6em]
  \item \textbf{Three of our bad-tip designs were wrong.} In the first, the planted tip stated the stale value, so agents submitted it without fetching anything and the ``spread'' was a one-hop trust test. In the second, the tip sat on a page that agents never opened (they list pages and read one), so it never spread, and the apparent run-to-run variation was luck about page names. In the third, the stale mirror warned about itself, so agents corrected. The final design uses a team-notes page everyone is asked to post to, a tip that is a link and never a value, a stale mirror with no warning, and an early-tip arm. We kept the run records of two of the discarded designs.
  \item \textbf{A strawman baseline.} Our first token tracer fell back to ``latest writer'' where a token was lost, while the baseline got the better rule. That made tokens look harmful for several models. The fair comparison gives both the same fallback, and the baseline is the best simple rule for each run.
  \item \textbf{A wrong answer key.} Defining the source as the most recent read of a route mislabelled agents that re-read pages (one model's lift looked like $+2$ points and is several times larger with the exact served copy). The key now prefers the exact served copy. Pooling worlds also hid regimes, so results are broken down by world.
  \item \textbf{A cherry-picked window.} Our first AI Village analysis used four hand-picked weeks and showed a post-rooms week at 87\% as typical; it was among the lowest. Running the whole dataset (Figure~\ref{fig:visibility}b) puts room-scoped coverage at \avPost{}. The claim about rooms is now a short dip, not a step change.
  \item \textbf{An early result that did not generalise.} A single-seed run with one Anthropic model (Haiku 4.5, 30 agents per arm) saw inert tags dropped in every copy and load-bearing tokens kept in every copy. The multi-seed, multi-family lab shows both depend on the model: the core open model kept inert tags in \qwenInert{} of copies. We no longer claim ``decoration fails and tokens work'' as a rule.
  \item \textbf{A hypothesis not met.} H1 (Table~\ref{tab:hyp}) failed, for the reason in Section~\ref{sec:res-survival}. H3 was not supported. H5 was dropped before analysis because the audit has nothing to say about swarms that share no distinctive strings.
  \item \textbf{A retracted claim.} We first wrote that the timing of a bad tip matters more than the model. With more families, an early tip took over only some of them (Figure~\ref{fig:tip}); the claim is withdrawn.
  \item \textbf{A tautological test, labelled as such.} The gated world traces by construction. We ran it to measure feasibility and cost and say so (Section~\ref{sec:res-gate}).
  \item \textbf{A model that could not take part.} One family never emitted a structured tool call in our harness and was dropped. Two others rarely copy, so they contribute little to the token results.
  \item \textbf{Rejected ideas.} We tested whether the swarm follows an ant-colony trail law (a herding exponent above one) and found an exponent of 1.3--1.5 without quality controls and 0.8 with per-method quality effects: the apparent herding was quality, not conformity. We also fitted a Hawkes process to adoption times in the incident, controlled for swarm-wide activity: adoption is largely self-triggered, and common outside triggers (a shared prompt, a scheduled start) cannot be excluded. It does not bear on the questions of this paper and we do not build on it.
\end{enumerate}
